\title{Unlocking Efficiency in Large Language Model Inference: Comprehensive Survey of Speculative Decoding}

\begin{abstract}
To mitigate the high inference latency stemming from autoregressive decoding in Large Language Models (LLMs), Speculative Decoding has emerged as a novel decoding paradigm for LLM inference. In each decoding step, this method first drafts several future tokens efficiently and then verifies them in parallel. Unlike autoregressive decoding, Speculative Decoding facilitates the simultaneous decoding of multiple tokens per step, thereby accelerating inference. This paper presents a comprehensive overview and analysis of this promising decoding paradigm. We begin by providing a formal definition and formulation of Speculative Decoding. Then, we organize in-depth discussions on its key facets, such as drafter selection and verification strategies. Furthermore, we present a comparative analysis of leading methods under third-party testing environments. We aim for this work to serve as a catalyst for further research on Speculative Decoding, ultimately contributing to more efficient LLM inference.\footnote{The relevant papers will be regularly updated at \url{https://github.com/hemingkx/SpeculativeDecodingPapers}.}
\end{abstract}

\section{Introduction}
Large Language Models (LLMs) have achieved remarkable proficiency in a range of downstream tasks~\cite{openai:2023gpt4, Hugo:2023llama, Hugo:2023llama2, vicuna2023, mistral}. They are progressively evolving as the cornerstone of comprehensive API interfaces (e.g., ChatGPT\footnote{\url{https://chat.openai.com}}), offering human life services and guidance through real-time human-machine interactions. However, the inference latency of these sizable models has emerged as a substantial obstacle restricting their broader applications. This latency primarily arises from the token-by-token generation necessitated by autoregressive decoding, resulting in an escalation of the inference latency with both the length of the generated sequence and the model's scale.

To accelerate LLM inference, an innovative inference paradigm, Speculative Decoding has been introduced~\cite{Stern:2018blockwise, xia:2022specdec, Leviathan:2023specdec, Chen:2023specsampling}. As shown in Figure~\ref{fig:specdec-intro}, in each decoding step, Speculative Decoding first efficiently drafts multiple tokens as speculation of future decoding steps of the target LLM and then utilizes the LLM to verify all drafted tokens in parallel. Only those tokens that meet the LLM's verification criterion are accepted as final outputs to guarantee generation quality. 

Speculative Decoding is founded upon two key observations about LLM inference: 1) many easy tokens can be predicted with less computational overhead (e.g., using a smaller model), and 2) LLM inference is highly memory bandwidth bound~\cite{Patterson:2004latencybandwith, Shazeer:2019memorybandwith} with the main latency bottleneck arising from memory reads/writes of LLM parameters rather than arithmetic computations. Drawing on these observations, Speculative Decoding adapts the concept of \textit{speculative execution}\footnote{Speculative execution ~\cite{Burton:1985specexe, John:2012specexe} is an optimization technique used in computer architecture where tasks are performed in advance and subsequently verified for their necessity, thereby circumventing the delays inherent in sequential task execution.} to focus LLMs' efforts on the validation of pre-drafted tokens, substantially diminishing the need for frequent memory operations of LLM parameters, thereby improving inference efficiency. 

While Speculative Decoding shows promise, it raises several critical questions that warrant further investigation. For instance, how to design an optimal drafter to strike a balance between speculation accuracy and drafting efficiency~\cite{xia:2022specdec, Zhou:2023distillspec, Li:2024eagle}. Additionally, it is essential to assess whether the verification criterion can maintain both generation parallelism and output quality~\cite{Miao:2023specinfer, medusa}. Furthermore, since existing methods are evaluated under disparate testing conditions, a unified benchmark is needed to facilitate realistic speedup expectations within the research community.

Amid the rapid expansion of research in Speculative Decoding, this work makes the first attempt to present a survey of this field, aiming to raise awareness among academics about the latest advancements. We provide a systematic categorization of current research and an in-depth analysis of relevant studies. Moreover, we introduce Spec-Bench, a comprehensive benchmark to assess Speculative Decoding methods in diverse application scenarios. Our contributions can be summarized as follows:

\begin{enumerate}[label={(\arabic*)}]
    \item \textbf{\textit{First survey:}} To our knowledge, we are the first to present a comprehensive survey on Speculative Decoding;
    \item \textbf{\textit{Formal definition:}} We furnish a formal definition and formulation of Speculative Decoding, laying the groundwork for future research.
    \item \textbf{\textit{New taxonomy:}} We provide a systematic taxonomy for Speculative Decoding, offering an organized categorization of existing work.
    \item \textbf{\textit{Spec-Bench:}} We introduce Spec-Bench, an extensive benchmark designed for assessing Speculative Decoding, enabling a comparative evaluation of leading methodologies.
\end{enumerate}

We hope that this work can serve as an essential guide for newcomers and motivate future research.

\section{Overview}
This paper offers a comprehensive survey of Speculative Decoding. We commence by introducing the early stages of Speculative Decoding research~(\S\ref{sec:timeline}), illustrated by a timeline of its evolution (as shown in Figure~\ref{fig:timeline}). This is followed by a formal definition and formulation of Speculative Decoding~(\S\ref{sec:definition}). Then, we delve into a detailed discussion of leading techniques, including the selection of draft models~(\S\ref{sec:drafting}), verification strategies~(\S\ref{sec:verification}), and alignment between the drafter and the target LLM~(\S\ref{sec:alignment}). Moreover, we introduce Spec-Bench, an extensive evaluation benchmark designed for assessing the acceleration effect of Speculative Decoding~(\S\ref{sec:comparisons}).

\section{Evolution of Speculative Decoding}
\label{sec:timeline}
This section discusses the motivation behind Speculative Decoding~(\S\ref{sec:motivation}) and then provides a detailed introduction to early attempts in this field~(\S\ref{sec:pioneering_efforts}).

\subsection{Motivation}
\label{sec:motivation}
The widespread adoption of LLMs has established autoregressive decoding as the \textit{de facto} standard to LLM inference~\cite{PaLM:2023, openai:2023gpt4, jiang2024mixtral}. However, autoregressive decoding is limited by its inference latency, which primarily stems from the memory-bound computation of LLMs~\cite{Patterson:2004latencybandwith, Shazeer:2019memorybandwith}. Specifically, the main latency bottleneck of each decoding step is not due to computational operations but arises from the necessity to transfer all LLM parameters from High-Bandwidth Memory (HBM) to the on-chip cache of modern accelerators like GPUs. This process, which generates only one token per step, leads to the underutilization of these accelerators and results in inefficiencies.

\subsection{Pioneering \textit{Draft-then-Verify} Efforts}
\label{sec:pioneering_efforts}
To mitigate the above issue, an intuitive way involves leveraging idle computational resources to enhance parallelism in LLM inference. To this end, \citet{Stern:2018blockwise} introduced Blockwise Decoding, an approach that incorporates extra feedforward neural (FFN) heads atop the Transformer decoder, enabling the simultaneous \textit{drafting} of multiple tokens per step. These tokens are then \textit{verified} by the original LLM \textit{in parallel}, ensuring that the outputs align with those of the original LLM. As a pioneering work proposing the \textit{Draft-then-Verify} paradigm, Blockwise Decoding effectively reduces the number of required LLM calls by increasing generation parallelism, thereby accelerating inference.

To further unleash the potential of this paradigm, \citet{xia:2022specdec} introduced Speculative Decoding (SpecDec), which utilizes an independent drafter, notably a specialized Non-Autoregressive Transformer, to perform the drafting task both accurately and efficiently. Moreover, this method presented an innovative strategy that relaxes the rigid verification criterion, thereby increasing the acceptance rate of drafted tokens. Impressively, SpecDec achieves around 5$\times$ speedup over autoregressive decoding with comparable quality, underscoring the substantial potential of Speculative Decoding.

\begin{algorithm}[t]
\small
\caption{Autoregressive Decoding}
\label{algo:ar}
\begin{algorithmic}[1]
\Require
Language model $\mathcal{M}_q$, input sequence $x_1,\dots,x_t$, and target sequence length $T$;
\State{\textbf{initialize} $n \gets t$}
\While{$n< T$}
\State Set $q_{n+1} \gets \mathcal{M}_q\left(x\mid x_{< n+1}\right)$ 
\State Sample $x_{n+1} \sim q_{n+1}$ 
\State $n \gets n + 1$
\EndWhile
\end{algorithmic}
\end{algorithm}

Following SpecDec, \citet{Leviathan:2023specdec} and \citet{Chen:2023specsampling} made concurrent contributions by proposing Speculative Sampling, expanding this paradigm to encompass the lossless acceleration of various sampling methods. These approaches employed smaller LMs from the same series (e.g., \texttt{T5-small}) to speed up the inference of their larger counterparts (e.g., \texttt{T5-XXL}). Unlike previous work, these \textit{off-the-shelf} small LMs do not require additional training, enabling the rapid adoption of Speculative Decoding in LLM acceleration. This advancement has elevated Speculative Decoding to the forefront of LLM efficiency research, attracting widespread interest within the NLP community.

To sum up, these pioneering efforts in Speculative Decoding have gradually solidified the \textit{Draft-then-Verify} paradigm, showcasing its promising potential in LLM acceleration. We provide a detailed categorization and discussion of these studies and subsequent research in the following sections.

\section{Formulation and Definition}
\label{sec:definition}
In this section, we first provide a concise overview of standard autoregressive decoding~(\S\ref{sec:ar}). Then, we offer an in-depth exposition of Speculative Decoding~(\S\ref{sec:specdec}), which encompasses a formal definition, a comprehensive description of the methodology, and a detailed elaboration of the algorithm.

\begin{algorithm}[t]
\small
\caption{Speculative Decoding}
\label{algo:specdec}
\begin{algorithmic}[1]
\Require
Target language model $\mathcal{M}_q$, draft model $\mathcal{M}_p$, input sequence $x_1, \dots, x_t$, block size $K$, target sequence length $T$, drafting strategy $\textsc{Draft}$, verification criterion $\textsc{Verify}$, and correction strategy $\textsc{Correct}$;
\State{\textbf{initialize} $n \gets t$}
\While{$n< T$}
\textcolor[RGB]{64,101,149}{\LineComment{Drafting: obtain distributions from $\mathcal{M}_p$ efficiently}}
\State Set $p_1, \dots, p_K \gets \textsc{Draft}\left(x_{\leq n}, \mathcal{M}_p\right)$
\textcolor[RGB]{64,101,149}{\LineComment{Drafting: sample $K$ drafted tokens}}
\State Sample $\widetilde{x}_i \sim p_i, i=1, \dots, K$
\textcolor[RGB]{64,101,149}{\LineComment{Verification: compute $K+1$ distributions in parallel}}
\State Set $q_{i} \gets \mathcal{M}_q\left(x\mid x_{\leq n}, \widetilde{x}_{<i}\right), i=1, \dots, K+1$
\textcolor[RGB]{64,101,149}{\LineComment{Verification: verify each drafted token}}
\For{$i=1:K$}
\If{$\textsc{Verify}\left(\widetilde{x}_i, p_{i}, q_{i}\right)$}
\State Set ${x}_{n+i} \gets\widetilde{x}_i$ and $n \gets n + 1$
\Else
\State ${x}_{n+i} \gets \textsc{Correct}\left(p_{i}, q_{i}\right)$
\State and Exit for loop.
\EndIf
\EndFor \State If all drafted tokens are accepted, sample next token ${x}_{n+1} \sim q_{K+1}$ and set $n \gets n + 1$.
\EndWhile
\end{algorithmic}
\end{algorithm}

\tikzstyle{my-box}=[
    rectangle,
    draw=hidden-draw, rounded corners, text opacity=1, minimum height=1.5em, minimum width=5em, inner sep=2pt, align=center, fill opacity=.5, line width=0.8pt, ]
\tikzstyle{leaf}=[my-box, minimum height=1.5em, fill=hidden-pink!80, text=black, align=left,font=\normalsize,
    inner xsep=2pt, inner ysep=4pt, line width=0.8pt, ]

\subsection{Autoregressive Decoding}
\label{sec:ar}
Transformer-based LLMs typically make generations in an autoregressive manner. Given an input sequence $x_1, \dots, x_t$, an autoregressive language model $\mathcal{M}_{q}$ generates the next token according to:

\begin{equation}\label{eq:ar}
    x_{t+1} \sim q_{t+1} = \mathcal{M}_q\left(x\mid x_{<t+1}\right),
\end{equation}

where $q$ is the conditional probability distribution calculated by $\mathcal{M}_{q}$ and $x_{t+1}$ denotes the next token sampled from $q_{t+1}$. We illustrate a detailed process in Algorithm~\ref{algo:ar}.

As discussed in Section~\ref{sec:timeline}, while the standard autoregressive decoding offers desirable generation quality, it is bounded by memory bandwidth, resulting in low utilization of modern accelerators. In this process, each memory-bound LLM call (i.e., an LLM forward step) produces merely a single token for the entire sequence, making the whole generation inefficient and time-consuming.

\subsection{Speculative Decoding}
\label{sec:specdec}
Following \citet{xia:2022specdec}, \citet{Leviathan:2023specdec}, and \citet{Chen:2023specsampling}, we here provide a formal definition of Speculative Decoding:

\begin{quote}
    Speculative Decoding is a \textit{Draft-then-Verify} decoding paradigm in which, at each decoding step, it first \textit{efficiently drafts} multiple future tokens and then \textit{verifies} all these tokens \textit{in parallel} using the target LLM to speed up inference.
\end{quote}

We formulate a detailed Speculative Decoding process in Algorithm~\ref{algo:specdec}. Subsequently, we delve into the two fundamental substeps integral to this paradigm -- \textit{drafting} and \textit{verification}:

\paragraph{Drafting} At each decoding step, Speculative Decoding first efficiently drafts multiple future tokens, as a speculation of the target LLM's output. Formally, given an input sequence $x_1, \dots, x_t$ and the target LLM $\mathcal{M}_{q}$, this paradigm employs an efficient draft model $\mathcal{M}_{p}$ (e.g., a smaller LM) to decode the next $K$ drafted tokens:

\begin{equation}
\begin{aligned}
p_1, \dots, p_K &= \textsc{Draft}\left(x_{\leq t}, \mathcal{M}_p\right), \\
\widetilde{x}_i &\sim p_i, \quad i=1, \dots, K,
\end{aligned}
\label{eq:draft-computing}
\end{equation}

where $\textsc{Draft}(\cdot)$ denotes various drafting strategies that we will discuss in Section~\ref{sec:drafting}, $p$ is the conditional probability distribution calculated by $\mathcal{M}_{p}$, and $\widetilde{x}_i$ denotes the drafted token sampled from $p_i$.

\paragraph{Verification} Subsequently, these drafted tokens are verified by the target LLM $\mathcal{M}_{q}$ in parallel. Formally, given the input sequence $x_1, \dots, x_t$ and the draft $\widetilde{x}_1, \dots, \widetilde{x}_K$, Speculative Decoding utilizes $\mathcal{M}_{q}$ to compute $K+1$ probability distributions simultaneously:

\begin{align}\label{eq:verify_computing}
q_{i} = \mathcal{M}_q\left(x\mid x_{\leq t}, \widetilde{x}_{<i}\right), i=1, \dots, K+1.
\end{align}

Then, each drafted token $\widetilde{x}_i$ is verified by a specific criterion $\textsc{Verify}\left(\widetilde{x}_i, p_{i}, q_{i}\right)$. Only those tokens that meet the criterion are selected as final outputs, ensuring quality consistent with the target LLM's standards. Otherwise, the first drafted token $\widetilde{x}_c$ that fails the verification will be corrected by the strategy $\textsc{Correct}\left(p_{c}, q_{c}\right)$. All drafted tokens after position $c$ will be discarded, to guarantee the high quality of the final outputs. If all tokens pass verification, an additional token $x_{t+K+1}$ will be sampled from $q_{K+1}$ as Eq.~(\ref{eq:ar}). 

The drafting and verification substeps will be iterated until the termination condition is met, i.e., the \texttt{[EOS]} token is decoded or the sentence reaches the maximal length. 

Notably, the acceleration effect of Speculative Decoding primarily hinges on \textit{the acceptance rate} of drafted tokens at each step. This rate is influenced by several factors, including the draft quality, verification criteria, and the behavior alignment between the drafter and the target LLM. Additionally, the intrinsic efficiency of the drafter itself also contributes to the overall speedup. In subsequent sections, we will delve into these pivotal components of Speculative Decoding, as depicted in Figure~\ref{fig:taxo_of_specdec}, to systematically categorize recent research trends within this promising paradigm.

\section{\texorpdfstring{}~Drafting}
\label{sec:drafting}
As a vital component of Speculative Decoding, the drafting process has a crucial impact on the speedup of the paradigm. The impact is determined by two key factors: the speculation accuracy of the drafter $\mathcal{M}_{p}$, measured by the average number of accepted tokens per step, and the drafting latency~\cite{Stern:2018blockwise, xia:2022specdec}. How to trade off high speculation accuracy and low drafting latency presents a major challenge in this process. In this section, we classify various drafting strategies into two categories: independent drafting~(\S\ref{sec:independent_drafting}) and self-drafting~(\S\ref{sec:self_drafting}), and summarize their formulations $\textsc{Draft}\left(x_{\leq t}, \mathcal{M}_p\right)$ in Table~\ref{tab:drafting}.

\subsection{Independent Drafting}
\label{sec:independent_drafting}
To strike a balance between speculation accuracy and efficiency, SpecDec~\cite{xia:2022specdec} first proposed utilizing an independent model for drafting. Specifically, it employed a specialized Non-Autoregressive Transformer that drafts multiple tokens simultaneously per step. This model has a deep-shallow encoder-decoder architecture to run efficiently. Despite its strengths, SpecDec requires training a draft model from scratch, which demands an increased computational budget.

Considering the available models in existing LLM series (e.g., OPT~\cite{opt} and LLaMA~\cite{Hugo:2023llama, Hugo:2023llama2}), a more straightforward and efficient approach is directly employing a small LM from the same series as the drafter to accelerate the inference of its larger counterparts~\cite{Leviathan:2023specdec, Chen:2023specsampling, Spector:2023stagedspec, Sun:2023spectr, chen2023cascadespec}. For instance, \citet{Leviathan:2023specdec} utilized \texttt{T5-small} as the drafter, to accelerate the inference of \texttt{T5-XXL}. These \textit{off-the-shelf} small LMs do not require additional training or any modification on model architectures, facilitating the quick adoption of Speculative Decoding. Moreover, since models in the same series share tokenizers, pretraining corpora, and similar training processes, they inherently have an alignment in prediction behaviors.

\subsection{Self-Drafting}
\label{sec:self_drafting}
While leveraging an external draft model offers considerable advantages, this approach necessitates extra effort to either train or identify a draft model that closely aligns with the target LLM. This challenge is intensified when no smaller counterparts of the LLM are available, e.g., LLaMA-7B~\cite{Hugo:2023llama, Hugo:2023llama2}. Furthermore, integrating two distinct models within a single system introduces additional computational complexity, particularly in distributed settings~\cite{medusa}. 

To address the above issues, numerous studies have suggested leveraging the target LLM itself for efficient drafting~\cite{Stern:2018blockwise, Santilli:2023paralleldecoding, Hooper:2023speed, medusa, Fu:2023lookahead, Du:2024glide}. Particularly, Blockwise Decoding~\cite{Stern:2018blockwise} and Medusa~\cite{medusa} incorporated FFN heads atop the Transformer decoder, enabling the parallel token generation per step. Compared with external drafters, these lightweight heads reduce extra computational overhead and are friendly to distributed inference. Another line of research has explored the potential of \textit{early exiting} and \textit{layer skipping} within the target LLM for drafting~\cite{Yang:2023PPD, Zhang:2023draftverify, Hooper:2023speed}. For instance, \citet{Yang:2023PPD} introduced additional subprocesses that exit early during the current decoding step, thereby initiating the drafting of future tokens in advance. Similarly, Self-Speculative~\cite{Zhang:2023draftverify} proposed to adaptively skip several intermediate layers during inference to draft efficiently. 

In contrast to prior work that focused on extending model architectures or altering the inference process, \citet{Santilli:2023paralleldecoding} introduced a simple drafting strategy that directly appends multiple \texttt{[PAD]} tokens to the end of the input prompt to enable parallel generation. However, this approach deviates from LLMs' autoregressive pretraining pattern, leading to suboptimal drafting quality. To tackle this, \citet{Fu:2023lookahead} proposed transforming low-quality drafts into multiple n-grams to improve the speculation accuracy; \citet{Monea:2023PaSS} introduced multiple learnable \texttt{[LA]} tokens and finetuned these token embeddings on a small training dataset to enhance the parallel decoding performance.

\section{\texorpdfstring{}~Verification}
\label{sec:verification}
In each decoding step, the drafted tokens are \textit{verified in parallel} to ensure the outputs align with the target LLM. This process also determines the number of tokens accepted per step, a vital factor impacting the speedup. This section summarizes various verification criteria $\textsc{Verify}\left(\widetilde{x}_i, p_i, q_i\right)$ (as shown in Table~\ref{tab:verification}), encompassing those supporting greedy decoding~(\S\ref{sec:greedy_decoding}) and speculative sampling~(\S\ref{sec:nucleus_sampling}) in LLM inference. Besides, we introduce token tree verification~(\S\ref{sec:token_tree_verification}), an effective strategy to increase the token acceptance rate.

\subsection{Greedy Decoding}
\label{sec:greedy_decoding}
Early attempts at Speculative Decoding focused on the verification criterion supporting greedy decoding, which guarantees that the outputs are exactly the same as the greedy decoding results of the target LLM~\cite{Stern:2019, Sun:2020SAD, xia:2022specdec}. Formally, given the input sequence $x_1, \dots, x_t$, the drafted tokens $\widetilde{x}_1, \dots, \widetilde{x}_K$, and the computed probability distributions $p_1, \dots, p_{K}$, $q_1, \dots, q_{K}$ as obtained from Eq.~(\ref{eq:draft-computing}) and (\ref{eq:verify_computing}), respectively, the verification criterion on the $i_{th}$ drafted token is formulated as

\begin{equation}\label{eq:verify-greedy}
\widetilde{x}_i = \arg \max q_i,
\end{equation}

where $i=1,\dots,K$. The first position $c$ that the drafted token $\widetilde{x}_c$ fails the verification denotes the \textit{bifurcation} position. The output token at this position ${x}_{t+c}$ will be adjusted by the correction strategy, which simply replaces the drafted token with the LLM's top-1 prediction:

\begin{equation}
{x}_{t+c} \gets \arg \max q_c.
\end{equation}

The verification criterion of greedy decoding is straightforward and clear. Thus, multiple subsequent studies have adopted this criterion to demonstrate the efficacy of their methodologies~\cite{Santilli:2023paralleldecoding, Yang:2023PPD, Hooper:2023speed, Zhang:2023draftverify, Fu:2023lookahead}. However, the strict matching requirement of this criterion often results in the rejection of high-quality drafted tokens, simply because they differ from the top-1 predictions of the target LLM, thereby constraining the speedup of the paradigm.

To tackle this problem, multiple studies have proposed various approximate verification criteria~\cite{Stern:2018blockwise, xia:2022specdec, Kim:2023bild}. Compared with the lossless criterion, these methods slightly relax the matching requirement to trust the drafts more, leading to higher acceptance of drafted tokens. For instance, SpecDec~\cite{xia:2022specdec} only requires the drafted tokens to fall in top-k candidates of the target LLM; BiLD~\cite{Kim:2023bild} proposed a rollback criterion that only rejects drafted tokens when the number of consecutive mismatch tokens exceeds a fixed threshold.

\subsection{Speculative Sampling}
\label{sec:nucleus_sampling}
Following \citet{Stern:2019}, subsequent work extended Speculative Decoding to support various sampling methods~\cite{Leviathan:2023specdec, Chen:2023specsampling}, accelerating the target LLM's inference without changing its output distribution. Formally, given the initial sequence $x_1, \dots, x_t$, the drafted tokens $\widetilde{x}_1, \dots, \widetilde{x}_K$ and the computed distributions $p_1, \dots, p_{K}$, $q_1, \dots, q_{K}$, the verification criterion on the $i_{th}$ drafted token is

\begin{equation}
\begin{aligned}\label{eq:verify-sampling}
r < \min\left(1, \frac{q_i(\widetilde{x}_i)}{p_i(\widetilde{x}_i)}\right), r \sim U\left[0,1\right],
\end{aligned}
\end{equation}

where $r$ denotes a random number drawn from a uniform distribution $U\left[0,1\right]$; $q_i(\widetilde{x}_i)$ and $p_i(\widetilde{x}_i)$ are the probability of $\widetilde{x}_i$ according to $\mathcal{M}_q$ and $\mathcal{M}_p$, respectively; and $i=1, \dots, K$. In other words, this criterion accepts the token $\widetilde{x}_i$ if $q_i(\widetilde{x}_i) \geq p_i(\widetilde{x}_i)$, and in case $q_i(\widetilde{x}_i) < p_i(\widetilde{x}_i)$ it rejects the token with probability $1-\frac{q_i(\widetilde{x}_i)}{p_i(\widetilde{x}_i)}$. The correction strategy resamples the output token at the bifurcation position $c$ from an adjusted distribution:

\begin{equation}
x_{t+c} \sim \operatorname{norm}(\max\left(0, q_c - p_c\right)).
\end{equation}

\citet{Leviathan:2023specdec} and \citet{Chen:2023specsampling} have theoretically proved that this criterion maintains identical output distributions to the target LLM. Thus, it has been widely adopted in subsequent research~\cite{Liu:2023onlinespec, Zhou:2023distillspec, Monea:2023PaSS, chen2023cascadespec}. In addition to the strict requirement, some work has also explored approximate strategies to improve the token acceptance rate~\cite{Leviathan:2023specdec, Zhou:2023distillspec}. For instance, \citet{Leviathan:2023specdec} proposed multiplying $p_i(\widetilde{x}_i)$ in Eq.~(\ref{eq:verify-sampling}) by a lenience parameter $l \in [0,1]$ to slightly relax the criterion.

\subsection{Token Tree Verification}
\label{sec:token_tree_verification}
Contrary to prior verification strategies that focused on a single draft sequence, SpecInfer~\cite{Miao:2023specinfer} proposed \textit{token tree verification}, an effective strategy enabling the target LLM to verify multiple draft sequences in parallel. As illustrated in Figure~\ref{fig:token_tree}, this method first merges multiple candidate draft sequences into a \textit{token tree} by sharing prefixes. It then utilizes a specially designed \textit{tree attention mask} to facilitate the LLM verifying the whole structure in parallel. Recent research has explored various approaches to obtain these candidate draft sequences~\cite{Miao:2023specinfer, medusa, He:2023REST, Li:2024eagle}. For instance, \citet{Miao:2023specinfer} generated diverse draft sequences from different boost-tuned LMs; \citet{medusa} considered the top-k predictions from each FFN head to obtain multiple candidate sequences.

\section{Alignment}
\label{sec:alignment}
As illustrated in Section~\ref{sec:drafting}, the speedup of Speculative Decoding primarily depends on the speculation accuracy, which in turn is influenced by the behavior similarity between the drafter and the target LLM. To enhance this, existing research has explored various knowledge distillation (KD) strategies to align the drafter's outputs with those of the target LLM~\cite{Stern:2018blockwise, xia:2022specdec, Miao:2023specinfer, Liu:2023onlinespec, Kim:2023bild, Zhou:2023distillspec}. Particularly, Blockwise Decoding adopted sequence-level knowledge distillation (Seq-KD)~\cite{kim2016seqkd} for alignment, which trained the drafter on the sentences generated by the target LLM. \citet{Miao:2023specinfer} proposed a collective boost-tuning (Col-BT) strategy, applying Seq-KD to finetune multiple small LMs on the training data and utilizing their aggregated output as drafts to improve the speculation accuracy.

Although Seq-KD is effective, it ignores the probability distributions of the target LLM, leading to performance degradation with sampling methods. To rectify this, recent studies have explored other KD strategies for Speculative Decoding~\cite{Zhou:2023distillspec, Liu:2023onlinespec}. Notably, DistillSpec~\cite{Zhou:2023distillspec} conducted a comprehensive comparison of different KD strategies on Speculative Decoding across various downstream tasks. \citet{Liu:2023onlinespec} proposed an online KD strategy that dynamically aligns the drafter with the target LLM on the fly using the query data.

We summarize the main features of existing Speculative Decoding methods in Table~\ref{tab:comparison}, including the drafter type or the drafting strategy, the alignment approach, supported verification strategies, and the reported speedup, etc. 

\section{Spec-Bench}
\label{sec:comparisons}

With the rapid research progress in Speculative Decoding, there is an increasing demand for comparative analysis of leading methods. However, existing approaches are tested using disparate benchmarks, devices, and environments, making fair comparisons impractical. To address this gap, we introduce Spec-Bench -- a comprehensive benchmark for Speculative Decoding covering diverse application scenarios. Based on Spec-Bench, we present a systematic comparison of open-source approaches under third-party testing conditions. Experiments were executed on \textit{the same device and testing environment} to ensure a fair comparison.

\subsection{Benchmark Construction} To assess Speculative Decoding methods across various scenarios, Spec-Bench encompasses six distinct subtasks: multi-turn conversation, translation, summarization, question answering, mathematical reasoning, and retrieval-augmented generation. We composed Spec-Bench by randomly selecting 80 instances from each of six widely used datasets, including MT-bench~\cite{zheng2023mtbench}, WMT14 DE-EN, CNN/Daily Mail~\cite{Nallapati:2016cnndm}, Natural Questions~\cite{kwiatkowski2019:naturalquestions}, GSM8K~\cite{Cobbe:2021gsm8k}, and DPR~\cite{Karpukhin:2020dpr}. For details on Spec-Bench and the specific experimental setup, please refer to Appendix~\ref{appendix:experimental_details}.

\subsection{Comparative Evaluation} Our main evaluations were conducted on Vicuna-7B at FP16 precision using a single consumer-grade 3090 GPU\footnote{For comparative analysis on a more powerful A100 GPU, please refer to Appendix~\ref{appendix:A100_analysis}.}. As depicted in Figure~\ref{fig:radar_3090}, under greedy settings, EAGLE~\cite{Li:2024eagle} achieves the highest speedup ratio (1.8$\times$$\sim$2.4$\times$) over autoregressive decoding across most subtasks, especially in mathematical reasoning (with a $\sim$2.4$\times$ speedup). EAGLE's success is mainly due to two factors: 1) it reuses the KV cache of LLMs to predict drafted tokens, substantially reducing the drafting computational overhead; 2) compared with Medusa~\cite{medusa}, EAGLE drafts in an autoregressive way, providing more stable and accurate speculation results. PLD~\cite{saxena2023pld} excels in subtasks with high similarities between input and output, such as summarization (with a $\sim$2.4$\times$ speedup). However, its performance diminishes in other subtasks like translation and question answering, with speedup ratios falling between 1.1$\times$$\sim$1.3$\times$.

We also compare the speedups of Speculative Decoding methods at different sampling temperatures. As illustrated in Figure~\ref{fig:sampling}, EAGLE consistently outperforms other methods across various settings, achieving a speedup ratio ranging from 1.7$\times$ to 2.1$\times$. Besides, it is observed that the acceleration effect of all methods decreases with an increase in sampling temperature. This is attributed to the increased computational complexity of the speculative sampling criterion at higher temperatures, as revealed in prior research~\cite{gante2023assisted, Spector:2023stagedspec}.

\section{Challenges and Future Directions}
\label{sec:directions}

\paragraph{How to trade off speculation accuracy and drafting efficiency?} As discussed in Sections~\ref{sec:drafting}, scaling up the drafter can effectively enhance speculation accuracy, yet it largely reduces the drafting efficiency and even the overall speedup. Therefore, it is essential to strike a balance between speculation accuracy and drafting latency. Among existing strategies, behavior alignment is a promising approach to address this issue, as it improves speculation accuracy without increasing latency. However, despite recent advancements~\cite{Miao:2023specinfer, Zhou:2023distillspec, Liu:2023onlinespec}, there is still considerable room for improvement to align the drafter with the target LLM. For example, given that the drafted tokens after the bifurcation position are all discarded, one potential direction could involve encouraging the drafter to prioritize the generation quality of early-position tokens. Beyond alignment, other factors such as the quality of drafting~\cite{Fu:2023lookahead} and the determination of speculation length~\cite{Su:2023batching} also influence speculation accuracy and merit further exploration.

\paragraph{How to apply Speculative Decoding in batched inference scenarios?} Currently, only a few Speculative Decoding implementations have supported batched inference, such as EAGLE\footnote{\url{https://github.com/SafeAILab/EAGLE}} and SpS\footnote{\url{https://github.com/lucidrains/speculative-decoding}}. However, batched inference is a crucial technique for efficiently managing user inputs in LLM real-time services. The primary challenges in batched Speculative Decoding lie in two aspects: (1) Each decoded sentence in Speculative Decoding varies in decoding steps due to different speculation accuracy. Thus, the inference latency of a batch depends on the slowest sample in the batch; (2) The extra computational complexity introduced by Speculative Decoding, especially in sampling settings, increases with larger batch sizes. How to maintain a promising speedup of Speculative Decoding in batched inference, and combine it with advanced techniques such as continuous batching~\cite{Yu:2022continuousbatching}, warrants further investigation.

\paragraph{How to integrate Speculative Decoding with other leading techniques?} As a general decoding paradigm, Speculative Decoding has already demonstrated its potential in conjunction with other advanced techniques~\cite{yang:2023llma, Zhang2023RaLMSpec, Li:2023contrastive}. For instance, \citet{Yuan:2023SCD} combined Speculative Decoding with Contrastive Decoding~\cite{Li:2023contrastive}, which not only speeds up the inference but also substantially improves the generation quality. In addition to the acceleration of text-only LLMs, applying Speculative Decoding in multimodal inference, such as image synthesis, text-to-speech synthesis, and video generation, is also an intriguing and valuable direction for future research. Another promising research direction is to integrate Speculative Decoding with other efficient methods such as vLLM~\cite{Kwon:2023vLLM}, Non-Auregressive Generation~\cite{Du:2021oaxe, Du:2022ngramxe} and Flash-Attention~\cite{Dao:2022flashattention, Dao:2023flashattention2}, further boosting the inference efficiency of LLM services.

\section{Conclusion}
\label{sec:conclusion}
This paper presents a comprehensive survey of Speculative Decoding, including the evolution of this promising paradigm, its formal definition and formulation, a systematic categorization of existing methods, and an in-depth review of leading techniques. Moreover, we introduce Spec-Bench, an extensive evaluation benchmark for Speculative Decoding methods, and present a comparative evaluation of prominent methods. To our knowledge, this is the first survey dedicated to Speculative Decoding. Our aim for this paper is to clarify the current research landscape and provide insights into future research directions.

\section*{Limitations}
\label{subsec:limitation} 

This paper provides a thorough examination and categorization of current methodologies and emerging trends in Speculative Decoding. We have also conducted a comparative analysis of leading open-source methods to offer researchers deeper insights into the advantages and limitations of different models. Beyond Speculative Decoding, we acknowledge additional efficient NLP strategies such as vLLM~\cite{Kwon:2023vLLM} and continuous batching~\cite{Yu:2022continuousbatching}. In the future, we intend to expand the discussion to encompass the integration of Speculative Decoding with these advanced techniques. Moreover, due to the absence of an available implementation of batched Speculative Decoding, our evaluations could not cover this aspect. We plan to undertake subsequent experiments to assess the speedup of Speculative Decoding methods across various batch sizes.

\section*{Ethics Statement}
\label{subsec:ethics} 
The datasets used in our experiment are publicly released and labeled through interaction with humans in English. In this process, user privacy is protected, and no personal information is contained in the dataset. The scientific artifacts that we used are available for research with permissive licenses. And the use of these artifacts in this paper is consistent with their intended use. Therefore, we believe that our research work meets the ethics of ACL. 

\end{document}
